% Content imported from Report_newest.docx.md with approved edits.
\chapter{CƠ SỞ LÝ THUYẾT VỀ KIẾN TRÚC ĐA THUÊ BAO}\label{v4-chapter-1}

Sự dịch chuyển của ngành phần mềm từ cài đặt tại chỗ sang cung ứng dịch vụ trên đám mây đã làm thay đổi cách thiết kế và vận hành các hệ thống thông tin doanh nghiệp. Bài toán đặt ra là làm thế nào để một hệ thống phần mềm duy nhất phục vụ đồng thời nhiều tổ chức với chi phí hợp lý, trong khi mỗi tổ chức vẫn có không gian dữ liệu, cấu hình nghiệp vụ và chất lượng phục vụ riêng. Kiến trúc đa thuê bao (Multi-Tenant Architecture) chính là lời giải cho bài toán này.

Chương~\ref{v4-chapter-1} trình bày cơ sở lý thuyết của kiến trúc đa thuê bao theo trình tự từ bài toán chi phí đến cấu trúc kỹ thuật. Cụ thể, chương giải thích lý do các dịch vụ SaaS cần chia sẻ tài nguyên (Mục~\ref{v4-sec-1-1}) rồi chuẩn hóa các khái niệm nền tảng để phân biệt bốn mô hình hệ thống và xác định mốc bắt đầu của đa thuê bao tại Cấp độ 3 (Mục~\ref{v4-sec-1-2}--1.3). Từ cơ sở đó, chương nêu năm đặc trưng cần có ở Cấp độ 3--4 (Mục~\ref{v4-sec-1-4}) và giới thiệu cấu trúc 9 tầng Cloud-Native cùng bảy thách thức kỹ thuật (Mục~\ref{v4-sec-1-5}), làm nền tảng cho việc lựa chọn kiến trúc ở Chương~\ref{v4-chapter-2}.

\section{Tổng quan SaaS và động lực chia sẻ tài nguyên}\label{v4-sec-1-1}

\subsection{Phần mềm dưới dạng dịch vụ và tiến trình chia sẻ tài nguyên}\label{v4-sec-1-1-1}

Trong lịch sử phát triển công nghệ phần mềm, phương thức cung cấp truyền thống dựa trên mô hình phần mềm đóng gói cài đặt tại chỗ (On-Premises). Ở mô hình này, mỗi tổ chức khách hàng phải mua giấy phép sử dụng theo phiên bản, tự trang bị hạ tầng máy chủ vật lý và duy trì đội ngũ kỹ thuật nội bộ để cài đặt, cấu hình cũng như bảo trì phần mềm trong phạm vi nội bộ. Điều này làm chi phí đầu tư ban đầu tăng cao, quy trình cập nhật phiên bản phân tán kéo dài, và sự lãng phí tài nguyên tính toán khi mỗi đơn vị phải dự phòng phần cứng theo mức tải cao điểm riêng lẻ.

Sự phổ cập của mạng Internet băng thông rộng kết hợp với công nghệ ảo hóa đã thúc đẩy phương thức cung ứng mới: Phần mềm dưới dạng Dịch vụ (Software as a Service - SaaS). SaaS là mô hình dịch vụ đám mây cho phép người sử dụng khai thác các ứng dụng của nhà cung cấp đang vận hành trên một hạ tầng điện toán đám mây thông qua giao diện mỏng như trình duyệt web hoặc giao diện lập trình ứng dụng (API). Trong đó khách hàng không trực tiếp quản lý hay kiểm soát hạ tầng mạng, máy chủ, hệ điều hành và phần lưu trữ bên dưới, ngoại trừ một số thiết lập cấu hình ứng dụng trong phạm vi được cấp quyền \cite{v4-14}. Nhờ chuyển đổi chi phí đầu tư vốn ban đầu sang chi phí thuê bao định kỳ và tập trung hóa công tác vận hành về phía nhà cung cấp, SaaS đã xóa bỏ rào cản tiếp cận công nghệ cho các doanh nghiệp.

Tuy nhiên, để nhà cung cấp SaaS có thể phục vụ số lượng lớn khách hàng mà không làm chi phí hạ tầng tăng tỷ lệ thuận với quy mô người dùng, hệ thống phần mềm cần có cơ chế chia sẻ tài nguyên hiệu quả. Dưới góc độ tiến hóa công nghệ, kiến trúc đa thuê bao hiện đại là kết quả kế thừa và khắc phục tuần tự những hạn chế của bốn giai đoạn chia sẻ tài nguyên \cite{v4-05}:

\begin{itemize}
\item
  \textbf{Thập niên 1960 (Chia sẻ thời gian trên máy tính lớn):} Nhiều thiết bị đầu cuối cùng chia sẻ thời gian xử lý của một CPU trung tâm. Tuy vậy, cách chia sẻ này mới dừng ở phạm vi máy tính nội bộ, chưa tạo thành mô hình cung cấp dịch vụ qua mạng diện rộng.
\item
  \textbf{Thập niên 1990 (Nhà cung cấp dịch vụ ứng dụng -- ASP):} Phần mềm được đặt tập trung tại trung tâm dữ liệu để khách hàng truy cập từ xa. Do kiến trúc ban đầu viết cho từng đơn vị riêng lẻ, nhà cung cấp vẫn phải cấp máy chủ, tiến trình và cơ sở dữ liệu riêng cho mỗi tổ chức, khiến chi phí tăng tuyến tính theo số lượng khách hàng.
\item
  \textbf{Thập niên 2000 (Ứng dụng Web đa người dùng):} Hệ thống phục vụ lượng lớn người dùng cá nhân đồng thời trên cùng một máy chủ web và cơ sở dữ liệu. Hạn chế của mô hình này là chỉ cung cấp chức năng cố định theo từng tài khoản cá nhân, thiếu không gian quản trị theo từng tổ chức và không hỗ trợ tùy biến quy trình cho doanh nghiệp.
\item
  \textbf{Hiện nay (SaaS đa thuê bao hiện đại):} Kết hợp việc chia sẻ môi trường chạy và cơ sở dữ liệu để giảm chi phí, đồng thời vẫn giữ được không gian làm việc cô lập, quản trị phân cấp và cấu hình riêng cho từng tổ chức khách hàng.
\end{itemize}

\VfourFigure{1.1}{assets/figures/v4/image2.png}{Tiến trình phát triển các mô hình chia sẻ tài nguyên điện toán (1960-nay).}{fig:v4-1-1}

Nhìn lại tiến trình trên, điểm thay đổi xuyên suốt là mức độ chia sẻ tài nguyên ngày càng sâu: từ chia sẻ phần cứng thô, gom nơi lưu trữ nhưng tách môi trường chạy (ASP), dùng chung môi trường cho người dùng cá nhân rời rạc, đến dùng chung nền tảng phần mềm có nhận biết từng tổ chức và cho phép tùy biến động.

\subsection{Động lực kinh tế - kỹ thuật}\label{v4-sec-1-1-2}

Để thấy rõ vì sao đa thuê bao giúp tối ưu chi phí cho SaaS, cần xem xét mối quan hệ giữa cách cấp phát tài nguyên và chi phí vận hành. Một dịch vụ SaaS vẫn có thể triển khai theo hướng cấp riêng môi trường độc lập cho từng khách hàng. Cách làm này giúp thiết kế đơn giản nhưng gặp ngay hai bất lợi về chi phí:

\begin{itemize}
\item
  \textbf{Mô hình chi phí của triển khai tách riêng:} Giả sử mỗi môi trường có cùng chi phí duy trì C\_\allowbreak{}single và chi phí chung có thể bỏ qua trong khoảng quy mô đang xét. Khi phục vụ N khách hàng trên N môi trường, chi phí minh họa là C(N) = N × C\_\allowbreak{}single.
\item
  \textbf{Lãng phí tài nguyên do dự phòng theo tải cao điểm:} Nếu mỗi khách hàng được cấp phát một lượng tài nguyên cố định đủ đáp ứng thời điểm hoạt động cao nhất của riêng họ, phần lớn năng lực tính toán đó sẽ không được khai thác trong những khoảng thời gian thông thường.
\end{itemize}

Ngược lại, khi nhiều khách hàng dùng chung một tập tài nguyên và thời điểm tải cao điểm của họ không trùng nhau, hệ thống có thể linh hoạt điều phối phần năng lực dư giữa các nhóm. Cách hoạt động này dựa trên nguyên lý ghép kênh thống kê (Statistical Multiplexing): tải tổng hợp được san phẳng, giúp phục vụ cùng một lượng khách hàng với ít phần cứng dự phòng hơn nhiều. Nhờ vậy, hàm chi phí có dạng C(N) = C\_\allowbreak{}fixed + δ × N, trong đó C\_\allowbreak{}fixed là chi phí đầu tư ban đầu cho kiến trúc đa thuê bao, còn δ là chi phí biên giả định trên đoạn hệ thống còn đủ năng lực phục vụ. Giá trị này có thể thay đổi khi phải bổ sung máy chủ hoặc nhân lực vận hành.

\VfourFigure{1.2}{assets/figures/v4/image3.png}{Tương quan tổng chi phí sở hữu (TCO) giữa mô hình đơn thuê bao và đa thuê bao.}{fig:v4-1-2}

Có thể thấy trên Hình~\ref{fig:v4-1-2}, nền tảng đa thuê bao đòi hỏi chi phí đầu tư ban đầu cao hơn để giải quyết các bài toán cô lập và tùy biến. Khi số lượng khách hàng còn ít, triển khai tách riêng từng hệ thống có thể rẻ hơn; nhưng khi vượt qua Điểm hòa vốn, mô hình đa thuê bao phát huy rõ lợi thế kinh tế theo quy mô. Lưu ý rằng Hình~\ref{fig:v4-1-2} là đồ thị định tính minh họa xu hướng: tổng chi phí đa thuê bao vẫn tăng chậm theo N, nhưng chi phí bình quân trên mỗi thuê bao (C\_\allowbreak{}MT(N)/N) tiến dần về δ trong mô hình minh họa với δ cố định. Hòa vốn và mức tiết kiệm phụ thuộc vào cấu hình, tải và chi phí thực tế.

\section{Các khái niệm nền tảng và định nghĩa đa thuê bao}\label{v4-sec-1-2}

Trong một hệ thống phần mềm dịch vụ, mối quan hệ giữa khách hàng, tính năng cung cấp và hạ tầng kỹ thuật xoay quanh ba khái niệm chính: Thuê bao, Dịch vụ và Thể hiện. Mặc dù các khái niệm này xuất hiện phổ biến, phạm vi của từng thuật ngữ đôi khi vẫn bị hiểu lẫn với nhau. Do đó, mục này làm rõ định nghĩa của cả ba khái niệm, đặc biệt là sự khác nhau giữa hai cấp độ của Thể hiện.

\subsection{Thuê bao (Tenant)}\label{v4-sec-1-2-1}

Trong các ứng dụng nhiều người dùng thông thường, tài khoản cá nhân cho biết ai đang thao tác. Còn trong hệ thống đa thuê bao, khái niệm Thuê bao đại diện cho một tổ chức khách hàng để xác định thao tác đó thuộc về đơn vị nào. Một tài khoản cá nhân có thể tham gia vào nhiều thuê bao khác nhau với vai trò độc lập. Thuê bao đóng vai trò là không gian tổ chức gom trọn nhóm thành viên, bộ cấu hình riêng, quy định phân quyền nội bộ và phần dữ liệu của đơn vị đó. Trong ứng dụng quản lý công việc Kanban của đề tài, mỗi không gian làm việc (workspace) của một tổ chức hay nhóm dự án là một thuê bao độc lập.

\subsection{Dịch vụ (Service)}\label{v4-sec-1-2-2}

Dịch vụ là tập hợp toàn bộ chức năng nghiệp vụ mà nhà cung cấp chuẩn hóa, đóng gói và đưa qua mạng cho các Thuê bao sử dụng, như dịch vụ quản lý công việc Kanban trong đề tài này. Khác với phần mềm viết riêng theo đơn đặt hàng cố định, dịch vụ SaaS được thiết kế hướng thị trường (sẽ trình bày cụ thể ở Chương~\ref{v4-chapter-4}) để vận hành tập trung trên hạ tầng dùng chung, kèm cam kết về tính sẵn sàng, an toàn dữ liệu và khả năng cấu hình theo từng tổ chức.

\subsection{Thể hiện (Instance)}\label{v4-sec-1-2-3}

Khi tiếp cận thuật ngữ \textbf{Instance (Thể hiện)}, người đọc thường hình dung đó là một bản ứng dụng đang nằm trên đĩa cứng và bộ nhớ RAM của máy chủ nhằm giúp phần mềm hoạt động. Tuy nhiên, trong phạm vi đề tài này, cách hiểu đó chỉ tương ứng với khái niệm \textbf{Thể hiện ứng dụng (Application Instance)}, tức là thể hiện của riêng thành phần ứng dụng đơn lẻ trong một tập hợp các thành phần. Để phân tích rõ các mô hình kiến trúc đa thuê bao, đề tài phân biệt tường minh giữa \textbf{thể hiện của cả một hệ thống đang vận hành} để cung cấp dịch vụ cho thuê bao và \textbf{thể hiện của từng thành phần kỹ thuật bên trong nó}:

\begin{enumerate}
\def\labelenumi{\arabic{enumi}.}
\item
  \textbf{Cấp độ vĩ mô -- Thể hiện Hệ thống (System Instance):} Chỉ toàn bộ một môi trường triển khai hoàn chỉnh phục vụ thuê bao, gộp từ phần cứng, cơ sở dữ liệu cho đến ứng dụng.
\item
  \textbf{Cấp độ vi mô -- Thể hiện Thành phần (Component Instance):} Chỉ từng thành phần kỹ thuật chạy độc lập ở mỗi tầng bên trong một Thể hiện Hệ thống, chẳng hạn Thể hiện ứng dụng hoặc Thể hiện cơ sở dữ liệu.
\end{enumerate}

Thuật ngữ \textbf{Thể hiện hệ thống} được dùng khi đề cập đến tổng thể môi trường phục vụ một hoặc nhiều thuê bao; trong khi tên cụ thể của từng tầng (như \textbf{Thể hiện ứng dụng}, \textbf{Thể hiện cơ sở dữ liệu}) được dùng khi phân tích việc tách hay gộp tài nguyên ở từng tầng kỹ thuật.

Phân biệt rõ hai cấp độ Thể hiện Hệ thống và Thể hiện Thành phần là cơ sở để Mục~\ref{v4-sec-1-3} tách bạch giữa Đa thể hiện (mỗi khách hàng một Thể hiện Hệ thống riêng) và Đa thuê bao (nhiều khách hàng dùng chung một Thể hiện Hệ thống). Đồng thời, cách tách này cũng làm tiền đề cho Chương~\ref{v4-chapter-2} phân tích phương án tách hay gộp từng Thể hiện Thành phần ở mỗi tầng kỹ thuật.

\subsection{Định nghĩa đa thuê bao và mô hình khái niệm}\label{v4-sec-1-2-4}

Tổng hợp góc nhìn của cả giới học thuật và công nghiệp, \textbf{Kabbedijk và cộng sự (2015) \cite{v4-05}} đã đúc kết định nghĩa chuẩn hóa về Đa thuê bao như sau:

\emph{``Multi-tenancy is a property of a system where multiple customers, so-called tenants, transparently share the system's resources, such as services, applications, databases, or hardware, with the aim of lowering costs, while still being able to exclusively configure the system to the needs of the tenant.''}

(Đa thuê bao là thuộc tính của một hệ thống trong đó nhiều khách hàng, được gọi là các thuê bao, chia sẻ một cách trong suốt các tài nguyên của hệ thống - chẳng hạn như dịch vụ, ứng dụng, cơ sở dữ liệu hoặc phần cứng - nhằm mục tiêu giảm chi phí, đồng thời vẫn cho phép cấu hình hệ thống một cách riêng biệt theo nhu cầu của từng thuê bao.)

Định nghĩa của Kabbedijk và cộng sự \cite{v4-05} chỉ ra bốn đặc điểm cốt lõi của một hệ thống đa thuê bao: (1) tính trong suốt đối với người dùng, (2) chia sẻ linh hoạt tài nguyên hệ thống, (3) giảm chi phí vận hành và (4) hỗ trợ cấu hình riêng theo nhu cầu từng thuê bao.

\VfourFigure{1.3}{assets/figures/v4/image4.png}{Mô hình khái niệm đa thuê bao và phạm vi Thể hiện Hệ thống.}{fig:v4-1-3}

Trên Hình~\ref{fig:v4-1-3}, mỗi Thuê bao hiện diện như một cột dọc bao gồm nhóm người dùng, ngữ cảnh cấu hình ở tầng ứng dụng và không gian dữ liệu ở tầng cơ sở dữ liệu. Tất cả các thuê bao này cùng đi xuyên qua các tầng trung gian như Middleware, hệ điều hành và lớp ảo hóa xuống tới phần cứng dùng chung của Thể hiện Hệ thống.

\section{Phân loại mô hình hệ thống và cấp độ trưởng thành SaaS}\label{v4-sec-1-3}

\subsection{Phân biệt bốn mô hình hệ thống (ST, MI, MU, MT)}\label{v4-sec-1-3-1}

Trong thực tế, khái niệm Đa thuê bao hay bị nhầm với ba mô hình gần gũi là Đơn thuê bao, Đa thể hiện và Đa người dùng. Nhìn từ bên ngoài, cả bốn mô hình đều cung cấp phần mềm qua mạng; nhưng nếu soi qua bộ ba khái niệm Thuê bao, Dịch vụ và Thể hiện ở Mục~\ref{v4-sec-1-2}, sự khác biệt bộc lộ rõ qua ba cặp dễ nhầm sau:

\begin{itemize}
\item
  \textbf{Nhầm lẫn giữa Đa người dùng và Đa thuê bao:} Cả hai đều cho phép nhiều người dùng chung một Thể hiện Hệ thống. Tuy nhiên, Đa người dùng không có khái niệm Thuê bao để gom nhóm theo tổ chức: mọi tài khoản cá nhân nằm phẳng trên cùng một không gian với chức năng cố định, thiếu cả tính cô lập cấp tổ chức lẫn khả năng tùy biến động dựa trên siêu dữ liệu.
\item
  \textbf{Nhầm lẫn giữa Đa thể hiện và Đa thuê bao:} Cả hai đều phục vụ nhiều tổ chức dưới cùng một thương hiệu Dịch vụ và cho phép cấu hình riêng. Nhưng mô hình Đa thể hiện cấp cho mỗi Thuê bao một Thể hiện Hệ thống tách rời, không gộp nhiều thuê bao vào cùng một môi trường nên hiệu quả sử dụng phần cứng bị hạn chế.
\item
  \textbf{Nhầm lẫn giữa Đơn thuê bao và Đa thuê bao:} Đơn thuê bao thực chất là đưa một phần mềm truyền thống lên máy chủ đám mây để thu phí định kỳ, chỉ thay đổi nơi đặt hạ tầng và cách tính phí chứ toàn bộ Thể hiện Hệ thống vẫn chỉ phục vụ duy nhất một tổ chức.
\end{itemize}

Từ ba cặp so sánh trên, có thể rút ra ba tiêu chí chính để phân biệt bốn mô hình:

\begin{enumerate}
\def\labelenumi{\arabic{enumi}.}
\item
  \textbf{Quản lý thuê bao và định danh hai lớp:} Hệ thống chỉ quản lý tài khoản cá nhân phẳng một lớp (user\_\allowbreak{}id như ở Đa người dùng), hay có thêm lớp tổ chức bên ngoài (tenant\_\allowbreak{}id bao bọc user\_\allowbreak{}id, cấu hình và dữ liệu như ở Đơn thuê bao, Đa thể hiện và Đa thuê bao)?
\item
  \textbf{Cách cấp phát Thể hiện Hệ thống:} Mỗi Thuê bao dùng riêng một Thể hiện Hệ thống (Đơn thuê bao, Đa thể hiện), hay nhiều Thuê bao dùng chung một Thể hiện Hệ thống thông qua phân giải ngữ cảnh động (Đa thuê bao)?
\item
  \textbf{Cách tùy biến hành vi nghiệp vụ:} Tùy biến bằng cách sửa mã nguồn hoặc cấu hình tĩnh theo từng bản cài đặt riêng (Đơn thuê bao, Đa thể hiện), không hỗ trợ tùy biến theo tổ chức (Đa người dùng), hay tùy biến động lúc chạy qua siêu dữ liệu trên một bộ mã nguồn duy nhất (Đa thuê bao)?
\end{enumerate}

\VfourFigure{1.4}{assets/figures/v4/image5.png}{So sánh bốn mô hình hệ thống (ST, MI, MU và MT).}{fig:v4-1-4}

Tổng hợp đối chiếu theo Hình~\ref{fig:v4-1-4} qua ba tiêu chí trên: Đơn thuê bao (ST) và Đa thể hiện (MI) đều có định danh tổ chức ở tiêu chí 1 nhưng phải cấp một Thể hiện Hệ thống riêng cho mỗi khách hàng ở tiêu chí 2, nên có thể cô lập tài nguyên tốt hơn, với chi phí xấp xỉ tỷ lệ số môi trường dưới giả định ở Mục~\ref{v4-sec-1-1-2}. Ngược lại, Đa người dùng (MU) gom mọi tài khoản vào một Thể hiện Hệ thống dùng chung nhưng thiếu ranh giới thuê bao hai lớp và khả năng tùy biến cho tổ chức. Chỉ kiến trúc Đa thuê bao (MT) dung hòa được cả ba tiêu chí: dùng chung một Thể hiện Hệ thống để giảm chi phí, giữ định danh hai lớp qua tenant\_\allowbreak{}id và tùy biến động theo siêu dữ liệu trên một bộ mã nguồn chung.

\subsection{Bốn cấp độ trưởng thành kiến trúc SaaS}\label{v4-sec-1-3-2}

\textbf{Chong và Carraro (2006) \cite{v4-15}} đã mô hình hóa lộ trình phát triển của kiến trúc phần mềm dịch vụ thành bốn cấp độ trưởng thành tuần tự dựa trên sự thay đổi của bộ mã nguồn, cơ chế tùy biến và cách thức tổ chức các thể hiện hệ thống:

\VfourFigure{1.5}{assets/figures/v4/image6.png}{Bốn cấp độ trưởng thành kiến trúc SaaS (Chong \& Carraro, 2006).}{fig:v4-1-5}

Quan sát Hình~\ref{fig:v4-1-5} và Bảng~\ref{tab:v4-1-1} có thể thấy sự chuyển dịch rõ nét qua bốn cấp độ. Ở Cấp độ 1 và Cấp độ 2, hệ thống vẫn cấp cho mỗi tổ chức một Thể hiện Hệ thống riêng (chỉ khác ở chỗ Cấp độ 2 gom về một bộ mã nguồn chuẩn với cấu hình tĩnh), nên chi phí vận hành phụ thuộc vào số môi trường phải duy trì. Bước chuyển lớn diễn ra tại Cấp độ 3 khi rào chắn phần cứng giữa các khách hàng được gỡ bỏ để gộp nhiều Thuê bao vào chung một Thể hiện Hệ thống, kết hợp tùy biến động bằng siêu dữ liệu lúc chạy. Cấp độ 4 kế thừa nguyên vẹn Cấp độ 3 để nhân bản thành nhiều thể hiện đặt sau bộ cân bằng tải. Như vậy, đa thuê bao đúng nghĩa bắt đầu từ Cấp độ 3. Cùng mạch nghiên cứu này \cite{v4-15}, Chong, Carraro và Wolter (2006) \cite{v4-16} cũng phân loại ba chiến lược tổ chức dữ liệu (cơ sở dữ liệu riêng, lược đồ riêng và bảng dùng chung), làm tiền đề cho phần phân tích ở Chương~\ref{v4-chapter-2}.

\begin{landscape}
\begingroup\VfourTableSetup
\renewcommand{\thetable}{1.1}
\begin{longtable}{@{}P{\dimexpr 0.20000\linewidth-8.000bp\relax}P{\dimexpr 0.20000\linewidth-8.000bp\relax}P{\dimexpr 0.20000\linewidth-8.000bp\relax}P{\dimexpr 0.20000\linewidth-8.000bp\relax}P{\dimexpr 0.20000\linewidth-8.000bp\relax}@{}}
\caption{So sánh bốn cấp độ trưởng thành kiến trúc SaaS (Chong \& Carraro, 2006).}\label{tab:v4-1-1}\\
\toprule
Tiêu chí Đối chiếu & Cấp độ 1: \emph{Tuỳ biến cục bộ} & Cấp độ 2: \emph{Cấu hình hoá} & Cấp độ 3: Đa thuê bao có khả năng cấu hình & \emph{Cấp độ 4: Đa thuê bao mở rộng quy mô} \\
\midrule\endfirsthead
\toprule
Tiêu chí Đối chiếu & Cấp độ 1: \emph{Tuỳ biến cục bộ} & Cấp độ 2: \emph{Cấu hình hoá} & Cấp độ 3: Đa thuê bao có khả năng cấu hình & \emph{Cấp độ 4: Đa thuê bao mở rộng quy mô} \\
\midrule\endhead
\midrule\endfoot
\bottomrule\endlastfoot
\textbf{1. Mô hình Hệ thống Tương ứng} & Đơn thuê bao tùy biến & Đa thể hiện. & \textbf{Đa thuê bao trên 1 thể hiện hệ thống.} & \textbf{\emph{Đa thuê bao mở rộng ngang (được phân tán / Cụm Đa thuê bao).}} \\
\textbf{2. Tổ chức Mã nguồn} & Nhiều nhánh mã nguồn riêng biệt theo từng khách hàng. & Một bộ mã nguồn chuẩn duy nhất & \textbf{Một bộ mã nguồn chuẩn duy nhất} & Một bộ mã nguồn chuẩn duy nhất \\
\textbf{3. Cấp phát Thể hiện (\emph{thể hiện hệ thống})} & \emph{Mỗi thuê bao 1 thể hiện hệ thống riêng biệt (môi trường).} & \emph{Mỗi thuê bao 1 thể hiện hệ thống riêng biệt (môi trường).} & \textbf{Nhiều thuê bao chia sẻ chung 1 \emph{thể hiện hệ thống} thống nhất.} & \emph{Cụm thể hiện hệ thống đa thuê bao sau bộ cân bằng tải.} \\
\textbf{4. Cơ chế Tùy biến \& Thuê bao} & Sửa trực tiếp mã nguồn; Thuê bao tách rời vật lý. & Cấu hình tĩnh theo bản cài đặt; Thuê bao tách rời theo \emph{thể hiện}. & \textbf{Tùy biến động lúc chạy qua Siêu dữ liệu; Thuê bao logic xuyên tầng.} & Kế thừa trọn vẹn cơ chế Siêu dữ liệu và logic của Cấp độ 3. \\
\textbf{5. Hiệu quả Chi phí (TCO) \& Mở rộng} & Chi phí vận hành và bảo trì cao nhất, tăng tuyến tính. & Giảm chi phí bảo trì mã nguồn, nhưng chi phí hạ tầng vẫn tăng. & \textbf{Tối ưu hóa mật độ phần cứng và TCO nhờ ghép kênh thống kê.} & Tối ưu hóa TCO kết hợp khả năng co giãn ngang lên quy mô cực lớn. \\
\end{longtable}
\endgroup
\end{landscape}

\section{Các đặc trưng cốt lõi của hệ thống đa thuê bao}\label{v4-sec-1-4}

Để vận hành an toàn và ổn định ở Cấp độ 3, đồng thời sẵn sàng mở rộng ngang lên Cấp độ 4, kiến trúc hệ thống cần đáp ứng đồng thời năm đặc trưng sau:

\subsection{Chia sẻ tài nguyên trong suốt và thực thi phi trạng thái}\label{v4-sec-1-4-1}

Chia sẻ tài nguyên trong suốt nghĩa là nhiều thuê bao cùng dùng chung một Thể hiện Hệ thống mà người dùng không cần biết hay can thiệp vào hạ tầng bên dưới. Hệ thống gộp luồng yêu cầu từ nhiều tổ chức lên cùng một tập tài nguyên: phía khách hàng thấy không gian làm việc riêng biệt, còn phía nhà cung cấp nâng cao được mức khai thác phần cứng nhờ bù trừ đỉnh tải giữa các tổ chức.

\textbf{Thực thi phi trạng thái để sẵn sàng lên Cấp độ 4:} Khi dùng chung một Thể hiện Hệ thống ở Cấp độ 3, các Thể hiện ứng dụng không được lưu trạng thái phiên làm việc của thuê bao trong RAM cục bộ của một máy chủ. Nhờ vậy, khi mở rộng thành nhiều máy chủ ở Cấp độ 4, mọi nút xử lý đều phục vụ được yêu cầu kế tiếp của bất kỳ thuê bao nào.

\subsection{Quản lý ngữ cảnh định danh và tùy biến động}\label{v4-sec-1-4-2}

\textbf{Nhận diện và giữ đúng mã thuê bao xuyên suốt:} Vì mọi yêu cầu từ nhiều tổ chức cùng đi vào một ứng dụng và một cơ sở dữ liệu dùng chung (khác với việc tách địa chỉ máy chủ riêng ở Cấp độ 1--2), hệ thống phải xác định chính xác tenant\_\allowbreak{}id ngay ở cổng vào và giữ nguyên mã này suốt vòng đời xử lý.

\textbf{Tùy biến động lúc chạy qua siêu dữ liệu:} Khác với Cấp độ 2 đọc cấu hình tĩnh lúc khởi động, một bộ mã nguồn ở Cấp độ 3 phải phục vụ xen kẽ nhiều tổ chức cùng lúc. Do đó, mọi thiết lập về giao diện, trường dữ liệu, quy tắc nghiệp vụ và hạn mức gói cước cần được đọc động theo tenant\_\allowbreak{}id từ kho siêu dữ liệu mà không phải sửa hay khởi động lại ứng dụng.

\subsection{Cô lập đa chiều (dữ liệu, hiệu năng, sự cố)}\label{v4-sec-1-4-3}

Khi gộp chung hạ tầng máy chủ và cơ sở dữ liệu vào một Thể hiện Hệ thống ở Cấp độ 3, việc bảo vệ ranh giới giữa các thuê bao chuyển vào trong phần mềm, đòi hỏi ba mức cô lập đi cùng nhau:

\begin{itemize}
\item
  \textbf{Cô lập dữ liệu và quyền truy cập:} Kiểm soát các thao tác đọc, thêm, sửa và xóa đúng theo tenant\_\allowbreak{}id, ngăn việc một tổ chức xem hoặc sửa nhầm dữ liệu của tổ chức khác dù nằm chung một Thể hiện Cơ sở Dữ liệu.
\item
  \textbf{Cô lập hiệu năng và tài nguyên:} Giới hạn mức dùng CPU, RAM, bộ nhớ đệm và bể kết nối cơ sở dữ liệu của từng thuê bao để tránh tình trạng ``hàng xóm ồn ào'' (Noisy Neighbor), tức một tổ chức gửi tải đột biến làm chậm tốc độ phản hồi của các tổ chức khác.
\item
  \textbf{Cô lập sự cố và khoanh vùng ảnh hưởng:} Tạo vách ngăn kỹ thuật để khi một thuê bao gặp lỗi xử lý hoặc quá tải, sự cố được giữ gọn trong phạm vi của thuê bao đó chứ không kéo sập toàn bộ hệ thống dùng chung.
\end{itemize}

\subsection{Tự phục vụ vòng đời và phân tách hai mặt phẳng}\label{v4-sec-1-4-4}

Để phục vụ nhiều tổ chức mà không cần can thiệp thủ công, hệ thống cần tự động hóa ba công việc vận hành: (1) cho phép đăng ký thuê bao mới và khởi tạo sẵn không gian dữ liệu cùng cấu hình gói cước; (2) co giãn linh hoạt số lượng tiến trình chạy theo lượng truy cập thực tế; và (3) nâng cấp phiên bản ứng dụng hoặc cấu trúc cơ sở dữ liệu tương thích ngược mà không làm gián đoạn các tổ chức đang sử dụng.

\textbf{Tách hai mặt phẳng vận hành (Control Plane và Data Plane):} Để dễ quản trị, cô lập sự cố và hỗ trợ tùy biến, toàn bộ chức năng phần mềm được chia thành hai mặt phẳng tách biệt:

\begin{enumerate}
\def\labelenumi{\arabic{enumi}.}
\item
  Mặt phẳng Điều khiển (Control Plane) gom các chức năng quản trị ở cấp nền tảng và cấp tổ chức như tạo hoặc tạm ngưng thuê bao, xác thực đăng nhập, quản lý gói cước, phân quyền thành viên và cài đặt cấu hình tùy biến. Mặt phẳng này phục vụ Quản trị viên Nền tảng (Platform Admin) và Quản trị viên Thuê bao (Tenant Admin), tần suất gọi vừa phải nhưng cần độ an toàn và toàn vẹn cao nhất.
\item
  Mặt phẳng Dữ liệu (Data Plane) đảm nhận các thao tác nghiệp vụ hằng ngày của người dùng cuối như mở bảng Kanban, cập nhật thẻ và kiểm tra quy tắc làm việc. Mặt phẳng này có tần suất truy cập lớn, cần độ trễ thấp và rào chắn dữ liệu chặt theo từng thuê bao. Việc tách hai mặt phẳng giúp tải cao ở Mặt phẳng Dữ liệu không làm gián đoạn công tác quản trị ở Mặt phẳng Điều khiển, còn mọi thay đổi cấu hình vẫn được đồng bộ tức thời xuống bên dưới mà không cần khởi động lại hệ thống.
\end{enumerate}

\subsection{Giám sát và đo lường theo thuê bao}\label{v4-sec-1-4-5}

Khi nhiều thuê bao dùng chung hạ tầng, việc chỉ đo các chỉ số chung của máy chủ (như tổng \% CPU hay dung lượng đĩa) sẽ che mất sự bất thường của từng tổ chức riêng lẻ. Hệ thống cần gắn thẻ mã thuê bao vào nhật ký và số liệu đo để phục vụ hai việc: (1) theo dõi thông lượng, độ trễ và tỷ lệ lỗi riêng của từng tổ chức; và (2) đối chiếu mức dùng thực tế với hạn mức gói cước để tính phí hoặc cảnh báo kịp thời.

\section{Mô hình phân tầng 9 tầng và các thách thức thiết kế}\label{v4-sec-1-5}

\subsection{Mô hình kiến trúc 9 tầng Cloud-Native}\label{v4-sec-1-5-1}

Thay vì dùng mô hình 3 tầng quen thuộc (Giao diện, Ứng dụng, Cơ sở dữ liệu) vốn quá gộp để phân biệt ranh giới tách hay dùng chung giữa máy chủ CSDL, CSDL logic, lược đồ và lớp trung gian, đề tài sử dụng khung 9 tầng của Kabbedijk và cộng sự (2015) \cite{v4-05} làm bản đồ tham chiếu chung. Ngoài ra, vì mô hình gốc ra đời trước khi công nghệ vùng chứa (Container) và các dịch vụ biên nhận biết thuê bao trở nên phổ biến, đề tài điều chỉnh Tầng 7 thành nét liền và bổ sung vùng chứa tại Tầng 3 để bám sát thực tế kiến trúc Cloud-Native hiện nay.

Trong công trình nghiên cứu năm 2015, \textbf{Kabbedijk và cộng sự \cite{v4-05}} đã giới thiệu mô hình phân tầng phần mềm gồm 9 tầng thành phần (Hình~\ref{fig:v4-1-6a}) nhằm cung cấp một khung tham chiếu chỉ ra các vị trí bên trong một thể hiện hệ thống mà tại đó cơ chế đa thuê bao có thể được áp dụng:

\VfourFigure{1.6a}{assets/figures/v4/image7.png}{Mô hình phân tầng 9 tầng nguyên bản của Kabbedijk và cộng sự (2015).}{fig:v4-1-6a}

Trong mô hình gốc (Hình~\ref{fig:v4-1-6a}), Kabbedijk và cộng sự chia hệ thống thành 9 tầng từ thấp đến cao. Các tác giả dùng nét liền cho 5 tầng phần mềm trực tiếp nhận biết đa thuê bao (Tầng 4, 5, 6, 8, 9) và nét đứt cho 4 tầng hạ tầng hoặc phần trung gian chưa đi sâu khi đó (Tầng 1, 2, 3 và Tầng 7).

Tuy nhiên, khi đặt vào bối cảnh kiến trúc hiện đại, mô hình nguyên bản cần được cập nhật tại hai điểm cốt lõi:

\begin{enumerate}
\def\labelenumi{\arabic{enumi}.}
\item
  Chuyển Tầng 7 từ nét đứt sang nét liền: Ở các hệ thống SaaS Cấp độ 3--4 hiện đại, các dịch vụ trung gian như cổng API (API Gateway), bộ nhớ đệm, hàng đợi thông điệp và kho lưu trữ tệp không chỉ truyền tải thụ động mà tham gia trực tiếp vào việc kiểm tra mã thuê bao, giới hạn tốc độ theo từng tổ chức và chia tách khóa đệm theo thuê bao.
\item
  \textbf{Bổ sung vùng chứa (Containers) tại Tầng 3 trên lớp máy ảo Tầng 2:} Bên cạnh máy ảo truyền thống, công nghệ vùng chứa đóng vai trò tiêu chuẩn đóng gói các dịch vụ trên cùng nhân hệ điều hành, hỗ trợ giới hạn phần cứng CPU và RAM rất gọn nhẹ.
\end{enumerate}

Trên cơ sở hai cập nhật trên, mô hình kiến trúc 9 tầng điện toán đám mây hiện đại được chuẩn hóa tại Hình~\ref{fig:v4-1-6b} và Bảng~\ref{tab:v4-1-2}:

\VfourFigure{1.6b}{assets/figures/v4/image8.png}{Mô hình kiến trúc đa thuê bao 9 tầng cập nhật cho Cloud-Native.}{fig:v4-1-6b}

\begin{landscape}
\begingroup\VfourTableSetup
\renewcommand{\thetable}{1.2}
\begin{longtable}{@{}P{\dimexpr 0.12000\linewidth-3.600bp\relax}P{\dimexpr 0.28000\linewidth-8.400bp\relax}P{\dimexpr 0.45000\linewidth-13.500bp\relax}P{\dimexpr 0.15000\linewidth-4.500bp\relax}@{}}
\caption{Vai trò của chín tầng thành phần trong kiến trúc đa thuê bao.}\label{tab:v4-1-2}\\
\toprule
Tầng thành phần & Bản chất kỹ thuật trừu tượng của Thể hiện Thành phần & Vai trò trong Kiến trúc Đa thuê bao Cấp độ 3 - 4 & Quy ước Phạm vi \\
\midrule\endfirsthead
\toprule
Tầng thành phần & Bản chất kỹ thuật trừu tượng của Thể hiện Thành phần & Vai trò trong Kiến trúc Đa thuê bao Cấp độ 3 - 4 & Quy ước Phạm vi \\
\midrule\endhead
\midrule\endfoot
\bottomrule\endlastfoot
\emph{\textbf{Tầng 9:} Thể hiện ứng dụng} & Tiến trình ứng dụng nghiệp vụ đang trực tiếp chạy để tiếp nhận và xử lý yêu cầu người dùng. & Ghép kênh nhiều thuê bao trên cùng tiến trình hoặc tách tiến trình; thực thi logic nghiệp vụ theo mã thuê bao. & \emph{\textbf{Nét liền} (Trọng tâm phần mềm)} \\
\emph{\textbf{Tầng 8:} Máy chủ ứng dụng} & Môi trường thực thi nền tảng quản lý luồng và vùng nhớ tiến trình. & Duy trì ngữ cảnh thuê bao cô lập trên từng luồng xử lý và nạp quy tắc cấu hình động. & \emph{\textbf{Nét liền} (Trọng tâm phần mềm)} \\
\emph{\textbf{Tầng 7:} Phần mềm trung gian} & Lớp dịch vụ điều phối trung gian: cổng giao tiếp lập trình, bộ nhớ đệm, hàng đợi thông điệp, kho lưu trữ đối tượng. & Trích xuất định danh thuê bao ở cửa ngõ, điều tiết lưu lượng và phân vùng khóa bộ nhớ đệm/hàng đợi/tệp. & \emph{\textbf{Nét liền} (Cập nhật điện toán đám mây hiện đại)} \\
\emph{\textbf{Tầng 6:} Lược đồ cơ sở dữ liệu} & Không gian định nghĩa cấu trúc bảng, thuộc tính, chỉ mục và chính sách ràng buộc dữ liệu. & Vị trí phân định giữa việc tách lược đồ riêng cho từng thuê bao và việc gộp chung bảng kèm chính sách lọc theo mã thuê bao. & \emph{\textbf{Nét liền} (Trọng tâm phần mềm)} \\
\emph{\textbf{Tầng 5:} Cơ sở dữ liệu logic} & Đơn vị cơ sở dữ liệu logic độc lập được khởi tạo bên trong một máy chủ CSDL. & Vị trí phân định giữa mô hình cấp riêng mỗi thuê bao một cơ sở dữ liệu và mô hình dùng chung cơ sở dữ liệu. & \emph{\textbf{Nét liền} (Trọng tâm phần mềm)} \\
\emph{\textbf{Tầng 4:} Máy chủ hệ quản trị cơ sở dữ liệu} & Tiến trình hệ quản trị cơ sở dữ liệu (như PostgreSQL) quản lý bộ đệm và kết nối. & Quyết định việc vận hành chung một cụm dịch vụ quản trị CSDL hay tách riêng tiến trình hệ quản trị cơ sở dữ liệu cho khách hàng lớn. & \emph{\textbf{Nét liền} (Trọng tâm phần mềm)} \\
\emph{\textbf{Tầng 3:} Hệ điều hành và vùng chứa} & Nhân hệ điều hành quản lý tiến trình và vào ra dữ liệu; hỗ trợ cô lập \emph{thể hiện vùng chứa}. & Cung cấp môi trường chạy vùng chứa/nhóm vùng chứa và kiểm soát hạn ngạch CPU và RAM cấp hệ điều hành. & \emph{\textbf{Nét đứt} (Hạ tầng nền)} \\
\emph{\textbf{Tầng 2:} Máy ảo} & Lớp ảo hóa phần cứng tạo ra các máy tính ảo độc lập. & Cung cấp ranh giới cô lập ở cấp máy ảo và mạng ảo trên hạ tầng đám mây. & \emph{\textbf{Nét đứt} (Hạ tầng nền)} \\
\emph{\textbf{Tầng 1:} Phần cứng vật lý} & Năng lực tính toán, bộ nhớ, lưu trữ và mạng vật lý thô. & Cung cấp nền tảng vật lý chịu lực dưới cùng cho toàn bộ các tầng phía trên. & \emph{\textbf{Nét đứt} (Hạ tầng nền)} \\
\end{longtable}
\endgroup
Để thuận tiện khi thiết kế giải pháp, 9 tầng thành phần trong Bảng~\ref{tab:v4-1-2} được gom thành bốn cụm theo màu sắc: Cụm Hạ tầng nền (Tầng 1--3, màu xám); Cụm Cơ sở dữ liệu (Tầng 4--6, màu xanh lá); Cụm Phần mềm trung gian (Tầng 7, màu cam); và Cụm Ứng dụng và Môi trường Thực thi (Tầng 8--9, màu xanh dương).
\end{landscape}

\subsection{Bảy thách thức thiết kế cốt lõi}\label{v4-sec-1-5-2}

Khi gộp chung các tầng ở Cấp độ 3 trên Hình~\ref{fig:v4-1-6b} và gỡ bỏ vách ngăn phần cứng giữa các tổ chức, hệ thống phải đối mặt với bảy thách thức kỹ thuật sau:

\begin{enumerate}
\def\labelenumi{\arabic{enumi}.}
\item
  \textbf{Thách thức 1 -- Lẫn lộn ngữ cảnh định danh qua các tầng:} Khi bể luồng ứng dụng ở Tầng 8 và bể kết nối CSDL ở Tầng 4 được luân phiên tái sử dụng cho nhiều thuê bao, chỉ cần thiếu bước kiểm tra ở Tầng 7 hoặc quên xóa ngữ cảnh sau khi xử lý xong là yêu cầu sau sẽ kế thừa nhầm tenant\_\allowbreak{}id của yêu cầu trước.
\item
  \textbf{Thách thức 2 -- Lộ dữ liệu chéo trên bảng và bộ đệm dùng chung:} Vì dữ liệu của mọi tổ chức nằm chung trong cùng các bảng ở Tầng 6 và cùng bộ nhớ đệm ở Tầng 7, lập trình viên chỉ cần viết sót một điều kiện lọc là dữ liệu có thể bị đọc chéo giữa các thuê bao (Mục~\ref{v4-sec-1-4-3}).
\item
  \textbf{Thách thức 3 -- Tranh chấp tài nguyên và hàng xóm ồn ào:} Một thuê bao gửi tải tăng vọt ở Tầng 7, chiếm quá nhiều bộ nhớ ở Tầng 8 hoặc chạy câu truy vấn quét rộng ở Tầng 6 sẽ trực tiếp làm chậm phản hồi của các thuê bao dùng chung bên cạnh.
\item
  \textbf{Thách thức 4 -- Giữ một bộ mã nguồn nhưng vẫn phải tùy biến lúc chạy:} Cần giữ duy nhất một bộ mã nguồn ở Tầng 9 và một lược đồ bảng ở Tầng 6 trong khi mỗi khách hàng lại cần giao diện, quy trình và trường dữ liệu riêng. Hệ thống không thể tách nhánh mã hay chạy lệnh sửa cấu trúc bảng vật lý trực tiếp vì sẽ khóa bảng và gây khó khăn khi bảo trì \cite{v4-17}.
\item
  \textbf{Thách thức 5 -- Tự động hóa vòng đời và đồng bộ hai mặt phẳng:} Không thể tạo thuê bao mới, đổi gói cước hay chỉnh cấu hình thủ công cho hàng trăm tổ chức. Mọi thao tác này phải tự động trên Mặt phẳng Điều khiển và đồng bộ nhanh xuống Mặt phẳng Dữ liệu (Mục~\ref{v4-sec-1-4-4}).
\item
  \textbf{Thách thức 6 -- Nâng cấp không gián đoạn và khôi phục dữ liệu riêng:} Khi nhiều tổ chức dùng chung một hệ thống đang chạy liên tục, mọi bản cập nhật ứng dụng ở Tầng 9 hay nâng cấp bảng ở Tầng 6 đều phải tương thích ngược để không làm gián đoạn dịch vụ \cite{v4-17}, đồng thời vẫn cần cách trích xuất hay khôi phục dữ liệu cho riêng một thuê bao khi cần.
\item
  \textbf{Thách thức 7 -- Khó quan sát tải và chi phí theo từng thuê bao:} Chỉ số CPU hay RAM tổng thể tại Tầng 1--4 không cho biết thuê bao nào đang gây nghẽn bên trong cụm ứng dụng dùng chung ở Tầng 7--9. Do đó, mọi nhật ký và chỉ số đo cần được gắn nhãn tenant\_\allowbreak{}id xuyên suốt quá trình chạy (Mục~\ref{v4-sec-1-4-5}).
\end{enumerate}

\section{Tiểu kết Chương~\ref{v4-chapter-1}}\label{v4-sec-1-6}

Chương~\ref{v4-chapter-1} đã trình bày trọn vẹn phần cơ sở lý thuyết từ tổng thể đến chi tiết: từ bài toán chi phí TCO và nguyên lý ghép kênh thống kê (Mục~\ref{v4-sec-1-1}), đến việc phân rõ Thuê bao -- Dịch vụ -- Thể hiện (Mục~\ref{v4-sec-1-2}), phân biệt bốn mô hình hệ thống cùng bốn cấp độ SaaS với mốc bắt đầu tại Cấp độ 3 (Mục~\ref{v4-sec-1-3}), năm đặc trưng của kiến trúc đa thuê bao (Mục~\ref{v4-sec-1-4}) và khung 9 tầng Cloud-Native kèm bảy thách thức kỹ thuật (Mục~\ref{v4-sec-1-5}).

Từ bảy thách thức này, câu hỏi đặt ra tiếp theo là: nên tách riêng hay gộp chung tài nguyên ở tầng nào để vừa tiết kiệm chi phí, vừa đáp ứng được bốn khía cạnh kiểm soát K1--K4, bài toán kiến trúc này sẽ được thảo luận ở Chương~\ref{v4-chapter-2}.
